% New section; it fits between the section on random graphs and the conclusion.
% It needs two macros next to \cp and \diam,
%   \newcommand{\tw}{\operatorname{tw}}
%   \newcommand{\torso}{\operatorname{torso}}
% and the entries Marx2013Separators and GareyJohnson1978Strong in references.bib.

\section{Approximation and Parameterized Complexity}
\label{sec:approx-param}

\subsection{Approximation Hardness}
\label{sec:approximation}

By Corollary~\ref{cor:two-approximation}, the union of a minimum $u$--$v$ cut and a shortest $u$--$v$ path is a cut-path with at most $2\cp(u,v) - 1$ edges, so a minimum cut-path is approximated within a factor of two in polynomial time.

A \emph{polynomial-time approximation scheme} (PTAS) for \MinCutPath{} is an algorithm that, given a graph $G = (V, E)$, two distinct vertices $u, v \in V$ and a number $\epsilon > 0$, returns a cut-path between $u$ and $v$ with at most $(1+\epsilon)\cp(u,v)$ edges, in time polynomial in $|V|$ for every fixed $\epsilon$.
It is a \emph{fully polynomial-time approximation scheme} (FPTAS) if its running time is polynomial in $|V|$ and $1/\epsilon$.

The only number in an instance of the decision version is the threshold $b$, and every instance with $b \geq |E|$ is a yes-instance, because $E$ is a cut-path.
% TODO(verify): the content of the cited paper (definition of strong NP-completeness, no FPTAS for such
% problems) was not read; only its bibliographic record was checked.
Hence the decision version is NP-complete in the strong sense of Garey and Johnson~\cite{GareyJohnson1978Strong}, and the standard argument for such problems excludes an FPTAS.

\begin{theorem}
\label{thm:no-fptas}
Unless $\mathrm{P} = \mathrm{NP}$, there is no FPTAS for \MinCutPath.
\end{theorem}

\begin{proof}
Suppose that an FPTAS exists.
Given an instance $(G, u, v, b)$ of the decision version with $G = (V, E)$, we run it on $(G, u, v)$ with $\epsilon = 1/(|E|+1)$ and obtain a cut-path $S$.
Since $\cp(u,v) \leq |E|$,
\[
\cp(u,v) \leq |S| \leq \Bigl(1 + \frac{1}{|E|+1}\Bigr)\cp(u,v) < \cp(u,v) + 1,
\]
and $|S| = \cp(u,v)$, because both numbers are integers.
The running time is polynomial in $|V|$ and $|E| + 1$, hence in the size of the instance.
Comparing $|S|$ with $b$ thus decides the decision version of \MinCutPath{} in polynomial time, and $\mathrm{P} = \mathrm{NP}$ by Theorem~\ref{thm:mcp-np-complete}.
\end{proof}

Theorem~\ref{thm:no-fptas} leaves open whether \MinCutPath{} has a PTAS and whether a polynomial-time algorithm achieves a factor smaller than two.
The reduction of Section~\ref{sec:np-completeness} gives no constant lower bound on the factor: in the graphs that it constructs, $\cp(u,v)$ is either $d(u,v)$ or at least $d(u,v) + 1$, and the ratio $(d(u,v) + 1)/d(u,v)$ tends to one as the formula grows.
Whether \MinCutPath{} is APX-hard, which would exclude a PTAS unless $\mathrm{P} = \mathrm{NP}$, is open.

\subsection{Fixed-Parameter Tractability}
\label{sec:fpt}

A decision problem is \emph{fixed-parameter tractable} (FPT) with respect to a parameter $b$ of its instances if it can be decided in time $f(b)$ times a polynomial in the size of the instance, for some computable function $f$ (see, e.g.,~\cite{Marx2013Separators}).
The \emph{parameterized version} of \MinCutPath{} is its decision version (Section~\ref{sec:ssp-to-mcp}) with the threshold $b$ as the parameter: given a graph $G = (V, E)$, two distinct vertices $u, v \in V$ and an integer $b$, is there a cut-path between $u$ and $v$ with at most $b$ edges?

By Lemma~\ref{lem:basic-bounds}, an instance with $d(u,v) > b$ or $c(u,v) > b$ is a no-instance, and an instance with $c(u,v) + d(u,v) - 1 \leq b$ is a yes-instance; both values are computed in polynomial time~\cite{Cormen2022}.

For a $u$--$v$ cut $C$, let $d_C(u,v)$ denote the minimum of $|P \setminus C|$ over all $u$--$v$ paths $P$; it is the distance between the images of $u$ and $v$ in the graph obtained from $G$ by contracting the edges of $C$.

\begin{lemma}
\label{lem:contraction}
Let $G = (V, E)$ be a graph, and let $u, v \in V$ be two distinct vertices. Then
\[
\cp(u,v) = \min_{C}\, \bigl(|C| + d_C(u,v)\bigr),
\]
where the minimum is over all inclusion-minimal $u$--$v$ cuts $C$.
\end{lemma}

\begin{proof}
For every $u$--$v$ cut $C$ and every $u$--$v$ path $P$, the set $C \cup P$ is a cut-path with $|C| + |P \setminus C|$ edges; hence $\cp(u,v) \leq |C| + d_C(u,v)$.
Conversely, a minimum cut-path $S$ contains a $u$--$v$ path $P$ and a $u$--$v$ cut, hence also an inclusion-minimal $u$--$v$ cut $C$, and
\[
|C| + d_C(u,v) \leq |C| + |P \setminus C| = |C \cup P| \leq |S| = \cp(u,v). \qedhere
\]
\end{proof}

By Lemma~\ref{lem:contraction}, one can decide whether $\cp(u,v) \leq b$ by computing $d_C(u,v)$ for every inclusion-minimal $u$--$v$ cut $C$ with at most $b$ edges, each time by a shortest-path computation in which the edges of $C$ have length zero.
This algorithm runs in time $|E|^{O(b)}$ and is not an FPT algorithm: a graph that consists of $b$ internally disjoint $u$--$v$ paths with $|E|/b$ edges each has $(|E|/b)^b$ inclusion-minimal $u$--$v$ cuts with $b$ edges.
We avoid the enumeration: after the remaining parts of the graph are replaced by weighted shortcuts, all these cuts lie in a graph of treewidth bounded by a function of $b$.

We use tree decompositions, treewidth and the monadic second-order logic of graphs as in~\cite[Section~2.1]{Marx2013Separators}, and write $\tw(H)$ for the treewidth of a graph $H$.
For a set $N$ of vertices of $H$, the graph $H - N$ is obtained from $H$ by deleting the vertices of $N$.
Let $x, y$ be two distinct vertices of $H$.
A set $N \subseteq V(H) \setminus \{x, y\}$ \emph{separates} $x$ from $y$ in $H$ if $x$ and $y$ lie in different components of $H - N$; such a set is an \emph{$x$--$y$ separator}, and it is \emph{minimal} if none of its proper subsets is an $x$--$y$ separator.
For $D \subseteq V(H)$, the graph $\torso(H, D)$ has the vertex set $D$, and two vertices $a, a' \in D$ are adjacent in it if they are adjacent in $H$ or if $H$ contains an $a$--$a'$ path whose internal vertices lie outside $D$~\cite[Definition~2.5]{Marx2013Separators}.

\begin{lemma}[{\cite[Proposition~2.7, Lemma~2.11 and Remark~2.13]{Marx2013Separators}}]
\label{lem:treewidth-reduction}
There are functions $f_1$ and $g_1$ and an algorithm with the following property.
Given a graph $H$, two distinct non-adjacent vertices $x, y$ of the same component of $H$, and an integer $k$ such that $H$ has an $x$--$y$ separator with at most $k$ vertices, the algorithm computes in time $f_1(k) \cdot (|V(H)| + |E(H)|)$ a set $D \subseteq V(H)$ with $x, y \in D$ such that
\begin{enumerate}[(a)]
    \item $D$ contains every minimal $x$--$y$ separator of $H$ with at most $k$ vertices,
    \item $\tw(\torso(H, D)) \leq g_1(k)$, and
    \item a set $N \subseteq D \setminus \{x, y\}$ separates $x$ from $y$ in $\torso(H, D)$ if and only if it separates $x$ from $y$ in $H$.
\end{enumerate}
\end{lemma}

Marx, O'Sullivan and Razgon~\cite[Lemma~2.11]{Marx2013Separators} construct the set $D \setminus \{x, y\}$ and bound the bramble number of its torso, which is the treewidth plus one.
Their bounds on the running time and on the bramble number depend on the minimum size of an $x$--$y$ separator, which is at most $k$, and on the difference between $k$ and this size; hence they are bounded by functions of $k$.
By their Remark~2.13, adding $x$ and $y$ increases the treewidth by at most two.
Statement~(c) is their Proposition~2.7.

\paragraph{Construction.}
Let $(G, u, v, b)$ be an instance with $c(u,v) \leq b$.
Let $G'$ be the graph obtained from $G$ by subdividing every edge $e \in E$ with a new vertex $z_e$, and let $R_F = \{z_e : e \in F\}$ for $F \subseteq E$.
Deleting the vertex $z_e$ from $G'$ corresponds to removing the edge $e$ from $G$.
Hence a set $F \subseteq E$ is a $u$--$v$ cut of $G$ if and only if $R_F$ separates $u$ from $v$ in $G'$, and $R_F$ is a minimal $u$--$v$ separator of $G'$ whenever $F$ is an inclusion-minimal $u$--$v$ cut of $G$.
The vertices $u$ and $v$ are not adjacent in $G'$ and lie in the same component, and $G'$ has a $u$--$v$ separator with $c(u,v) \leq b$ vertices.
We apply Lemma~\ref{lem:treewidth-reduction} to $G'$ with $x = u$, $y = v$ and $k = b$ and obtain a set $D$.
By~(a), $z_e \in D$ for every edge $e$ of an inclusion-minimal $u$--$v$ cut of $G$ with at most $b$ edges.

Let $G_D = \torso(G', D)$.
For an edge $aa'$ of $G_D$, let $\omega(aa')$ be the minimum number of vertices of $R_E$ among the internal vertices of an $a$--$a'$ path of $G'$ whose internal vertices lie outside $D$; in particular, $\omega(aa') = 0$ if $a$ and $a'$ are adjacent in $G'$.
Let $G_D^+$ be the graph obtained from $G_D$ by replacing every edge $aa'$ with $\omega(aa') \geq 1$ by an $a$--$a'$ path with $\min\{\omega(aa'), b+1\}$ new internal vertices.
Let $Y$ be the set of all new vertices, and let $B = R_E \cap D$.
Subdividing edges does not change which vertices of $D$ lie in a common component after vertices of $D$ are deleted.
Together with~(c) and the correspondence between the cuts of $G$ and the separators of $G'$, this gives, for every $F \subseteq E$ with $R_F \subseteq D$,
\begin{equation}
\label{eq:cut-separator}
\text{$F$ is a $u$--$v$ cut of $G$ if and only if $R_F$ separates $u$ from $v$ in $G_D^+$.}
\end{equation}
A set $M \subseteq B \cup Y$ is \emph{good} if
\begin{enumerate}[(i)]
    \item $M \cap B$ separates $u$ from $v$ in $G_D^+$, and
    \item $G_D^+$ contains a $u$--$v$ path all of whose vertices in $B \cup Y$ belong to $M$.
\end{enumerate}

\begin{lemma}
\label{lem:good-set}
Let $(G, u, v, b)$ be an instance with $c(u,v) \leq b$, and let $G_D^+$, $B$ and $Y$ be constructed as above.
Then $\cp(u,v) \leq b$ if and only if there is a good set with at most $b$ vertices.
\end{lemma}

\begin{proof}
Suppose first that $\cp(u,v) \leq b$.
By Lemma~\ref{lem:contraction}, there are an inclusion-minimal $u$--$v$ cut $C$ and a $u$--$v$ path $P$ of $G$ with $|C| + |P \setminus C| \leq b$.
Since $|C| \leq b$, we have $R_C \subseteq D$ and hence $R_C \subseteq B$.
Let $P'$ be the $u$--$v$ path of $G'$ obtained from $P$ by subdividing its edges, and let $u = a_0, a_1, \dots, a_r = v$ be the vertices of $P'$ that lie in $D$, in their order on $P'$.
For $1 \leq i \leq r$, the part of $P'$ between $a_{i-1}$ and $a_i$ has no internal vertex in $D$.
Hence $a_{i-1}a_i$ is an edge of $G_D$, and $\omega(a_{i-1}a_i)$ is at most the number of internal vertices of this part that lie in $R_E$.
Every such vertex is $z_e$ for an edge $e \in P$ with $z_e \notin D$, and then $e \notin C$.
In particular, $\omega(a_{i-1}a_i) \leq |P \setminus C| \leq b$, so $G_D^+$ contains a $u$--$v$ path $Q$ that passes through $a_0, \dots, a_r$ in this order and has exactly $\omega(a_{i-1}a_i)$ vertices of $Y$ between $a_{i-1}$ and $a_i$.
Let $M = R_C \cup \bigl(V(Q) \cap (B \cup Y)\bigr)$.
Condition~(ii) holds by the choice of $M$.
Condition~(i) holds because $M \cap B \supseteq R_C$ and $R_C$ separates $u$ from $v$ in $G_D^+$ by~\eqref{eq:cut-separator}.
The vertices of $Q$ in $B$ are the vertices $z_e \in D$ with $e \in P$, the number of vertices of $Q$ in $Y$ is at most the number of edges $e \in P$ with $z_e \notin D$, and $z_e \in D$ for every $e \in C$; hence
\[
|M| \leq |C| + |\{e \in P \setminus C : z_e \in D\}| + |\{e \in P \setminus C : z_e \notin D\}| = |C| + |P \setminus C| \leq b.
\]

Conversely, let $M$ be a good set with $|M| \leq b$, and let $Q$ be a $u$--$v$ path of $G_D^+$ as in~(ii).
Let $F_1 = \{e \in E : z_e \in M \cap B\}$; by~(i) and~\eqref{eq:cut-separator}, $F_1$ is a $u$--$v$ cut of $G$.
The path $Q$ consists of edges $aa'$ of $G_D$ with $\omega(aa') = 0$ and of paths that replace edges $aa'$ of $G_D$ with $\omega(aa') \geq 1$.
The internal vertices of such a path lie in $Y$, hence in $M$ by~(ii); since $|M| \leq b$, their number is not $b+1$, so it is $\omega(aa')$.
In both cases we replace the part of $Q$ between $a$ and $a'$ by an $a$--$a'$ path of $G'$ whose internal vertices lie outside $D$ and exactly $\omega(aa')$ of them in $R_E$.
The result is a $u$--$v$ walk $Q'$ in $G'$.
Let $F_2 = \{e \in E : z_e \in V(Q')\}$.
A vertex $z_e$ of $Q'$ is a vertex of $Q$ in $B$, and then $e \in F_1$ by~(ii), or it is an internal vertex of one of the inserted paths, and the number of these vertices is at most $|V(Q) \cap Y| \leq |M \cap Y|$.
Hence $|F_1 \cup F_2| \leq |M \cap B| + |M \cap Y| = |M| \leq b$.
In $G'$, the vertex $z_e$ is adjacent exactly to the two ends of $e$; therefore any two vertices of $V$ that follow each other on $Q'$ are equal or are the two ends of an edge of $F_2$, and $F_2$ contains a $u$--$v$ path.
Thus $F_1 \cup F_2$ is a cut-path between $u$ and $v$ with at most $b$ edges.
\end{proof}

\begin{theorem}
\label{thm:fpt}
\MinCutPath{} is fixed-parameter tractable with respect to the threshold $b$.
\end{theorem}

\begin{proof}
The algorithm computes $c(u,v)$ and rejects if $c(u,v) > b$, which is correct by Lemma~\ref{lem:basic-bounds}.
Otherwise it constructs $G'$, $D$, $G_D$, $\omega$ and $G_D^+$ and accepts if and only if there is a good set with at most $b$ vertices, which is correct by Lemma~\ref{lem:good-set}.

The graph $G'$ has $|V| + |E|$ vertices and $2|E|$ edges, so $D$ is computed in time $f_1(b) \cdot (|V| + 3|E|)$.
For two vertices $a, a' \in D$, a shortest-path computation in $G' - (D \setminus \{a, a'\})$, in which the length of a path is the number of its internal vertices in $R_E$, decides whether $aa'$ is an edge of $G_D$ and gives $\omega(aa')$.
Hence $G_D^+$ is constructed in polynomial time and has at most $|D| + (b+1)|D|^2$ vertices.

Subdividing an edge $aa'$ of a graph $H$ by a vertex $w$ yields a graph of treewidth at most $\max\{\tw(H), 2\}$.
Indeed, a tree decomposition of $H$ has a bag that contains $a$ and $a'$, and a new bag $\{a, a', w\}$ attached to it gives a tree decomposition of the new graph.
Hence $\tw(G_D^+) \leq \max\{\tw(G_D), 2\} \leq \max\{g_1(b), 2\}$ by~(b).

Let $W$ be a set of vertices of $G_D^+$ with $v \in W$.
The vertex $v$ is reachable from $u$ by a path whose vertices other than $u$ lie in $W$ if and only if every set of vertices that contains $u$ and every vertex of $W$ adjacent to one of its members also contains $v$.
This condition is a formula of monadic second-order logic.
Condition~(ii) is this condition for $W = M \cup \bigl(V(G_D^+) \setminus (B \cup Y)\bigr)$, condition~(i) is its negation for $W = V(G_D^+) \setminus (M \cap B)$, and $|M| \leq b$ holds if and only if there are $b$ vertices, not necessarily distinct, such that every member of $M$ is one of them.
Hence the existence of a good set with at most $b$ vertices is expressed by a formula $\varphi_b$ of monadic second-order logic over graphs with the vertex labels $u$, $v$, $B$ and $Y$.
By Courcelle's theorem~\cite[Theorem~2.2]{Marx2013Separators} and its extension to graphs with vertex labels stated there, a property expressed by such a formula $\varphi$ is decided on a graph $H$ in time $f_\varphi(\tw(H)) \cdot (|V(H)| + |E(H)|)$.
The total running time is therefore $f(b)$ times a polynomial in $|V| + |E|$, for a function $f$ of $b$.
\end{proof}

The proof gives no explicit bound on $f$.
Whether \MinCutPath{} can be solved in time $2^{O(b)}$ times a polynomial in $|V| + |E|$ is open.
